% GPT-6 Astra Pro; effort not exposed; /root/replacement; 22 November 2026.
\section[HD-F1-R34: the rank-three/rank-two gate is excluded]{HD-F1-R34: the remaining three-map adjacent gate is excluded}
\label{sec:followup-rm-three-forward-maps-exclusion}

There is no actual object in the
rank-$(3,2)$, equal-slope semistable TF
adjacent gate of R21--R28, with the
specified coherent faithful finite
point action and weakly semiregular
gluable lift. R33 excluded its
two-map alternatives. We now exclude
the remaining common-image case
\[
                       h=\dim\operatorname{Hom}_G(A,T)=3.
 \tag{F1R34.1}
\]
All actual source and target defects,
the original connecting morphism and
the full balanced character are
retained. The argument does not
address higher endpoint ranks,
other cohomological amplitudes,
or endpoints outside this equal-slope
semistable TF gate.

The improved full-character estimate
is twenty-eight. The exact triangle
maps can lose at most seven
dimensions at $h=3$, so they would
force
\[
                    \operatorname{ext}^2_G(E,E)\ge28-7=21,
 \tag{F1R34.2}
\]
contradicting the balanced value
eighteen. The loss of seven, including
all possible endpoint Yoneda kernels,
is proved below.

\paragraph{1. The actual mixed-character pencils.}
Keep the original triangle
\[
 A\xrightarrow{i}E\xrightarrow{p}C=T[-1]
       \xrightarrow{\delta[-1]}A[1],\qquad
                 \delta:T\longrightarrow A[2].
 \tag{F1R34.3}
\]
The determinant twist, R26's full
character classification and R27--R28's
actual cycle obstruction give
\[
 \begin{split}
 z_T&=a_0v_1v_2+c p_1p_2,\\
 z_A&=(a_0+1)v_1v_2+c p_1p_2,
 \end{split}
 \qquad a_0>0,\quad c\ne0,\quad
                p_i\notin\mathbf C v_i.
 \tag{F1R34.4}
\]
Use R33's actual exterior module
maps and notation
$\xi_i=\Phi_i^{-1}(v_i)$,
$\alpha_i=\Phi_i^{-1}(p_i)$
on the four-dimensional spaces $U_i$.
For a two-vector $\gamma\in\bigwedge^2U_i$
put
\[
 L_\gamma:U_i\longrightarrow\bigwedge^3U_i,
                     \quad u\longmapsto u\gamma.
 \tag{F1R34.5}
\]
The explicit balanced form
$\xi_i=-t_{i1}t_{i2}-q w_{i1}w_{i2}$
makes $L_{\xi_i}$ an isomorphism,
as checked on all four basis vectors
in R33. Define actual endomorphisms
of these cohomological source spaces by
\[
                  P_i=L_{\xi_i}^{-1}L_{\alpha_i},
                       \qquad \mathcal P=P_1\otimes P_2.
 \tag{F1R34.6}
\]
The map from two-vectors to the
operators (F1R34.5) is injective:
wedging with the four basis vectors
detects every coefficient of a
two-vector. Hence $P_i$ is not
scalar, because $\alpha_i$ and
$\xi_i$ are independent.

Postcomposing the two mixed-block
character maps with the same actual
inverse of $L_{\xi_1}\otimes L_{\xi_2}$
identifies them with
\[
                        a_0 I+c\mathcal P,
                   \qquad (a_0+1)I+c\mathcal P.
 \tag{F1R34.7}
\]
Thus their nullities are the dimensions
of the two actual eigenspaces of
$\mathcal P$ at
\[
                  u=-\frac{a_0}{c},\qquad
                  v=-\frac{a_0+1}{c}.
 \tag{F1R34.8}
\]
These numbers are nonzero and
distinct, and $u\ne-v$, since
$a_0>0$. No factor in (F1R34.4)
has been substituted for a sheaf;
(F1R34.6)--(F1R34.7) are actual
linear isomorphisms on the original
character-map spaces.

\paragraph{2. The complete four-dimensional pencil calculation.}
Fix a volume form on $U_i$ and
define the nondegenerate alternating
form
\[
                      \Omega_i(x,y)=[xy\xi_i].
 \tag{F1R34.9}
\]
Equation (F1R34.6) gives
\[
 \Omega_i(P_i x,y)=[xy\alpha_i]
                           =\Omega_i(x,P_i y).
 \tag{F1R34.10}
\]
The signs follow by moving the
two-vector past one vector, which
has positive sign. Thus $P_i$ is
self-adjoint for this alternating
form.

Choose a symplectic basis for
$\Omega_i$: select a vector and
a partner of nonzero pairing,
rescale that pairing to one, and
repeat on their two-dimensional
orthogonal complement. In the
resulting basis, with
\[
 J_4=\begin{pmatrix}0&I_2\\-I_2&0\end{pmatrix},
                  \qquad J_2=\begin{pmatrix}0&1\\-1&0\end{pmatrix},
\]
the equality $P_i^tJ_4=J_4P_i$
has exactly the block form
\[
                 P_i=\begin{pmatrix}B_i&b_iJ_2\\d_iJ_2&B_i^t\end{pmatrix}
                       \quad(B_i\in\operatorname{Mat}_2(\mathbf C)).
 \tag{F1R34.11}
\]
This follows by multiplying the two
block matrices: the off-diagonal
blocks must be alternating, and
the lower diagonal block is the
transpose of the upper one.

The two-by-two identities
\[
 \begin{gathered}
 B_i^2-\operatorname{tr}(B_i)B_i+\det(B_i)I_2=0,\\
 B_iJ_2+J_2B_i^t=\operatorname{tr}(B_i)J_2,
 \qquad J_2^2=-I_2
 \end{gathered}
\]
now give the full operator identity
\[
          P_i^2-\operatorname{tr}(B_i)P_i
                  +(\det(B_i)+b_i d_i)I_4=0.
 \tag{F1R34.12}
\]
For example, the upper diagonal
block of $P_i^2$ is $B_i^2-b_i d_iI_2$;
the off-diagonal blocks are
$\operatorname{tr}(B_i)$ times
their counterparts in $P_i$.
This proves (F1R34.12) over
$\mathbf C$ directly.

Because $P_i$ is not scalar, its
minimal polynomial has degree two.
There are exactly the following
two possibilities.

If its two roots $\lambda_i,\mu_i$
are distinct, the actual projectors
\[
     \frac{P_i-\mu_i I}{\lambda_i-\mu_i},\qquad
     \frac{P_i-\lambda_i I}{\mu_i-\lambda_i}
 \tag{F1R34.13}
\]
give the two eigenspace summands.
They are nonzero. Distinct
eigenspaces are orthogonal for
$\Omega_i$ by (F1R34.10), so the
restriction of $\Omega_i$ to each
is nondegenerate. Each therefore
has even dimension, and both
have dimension two.

If the root is repeated, write
$P_i=\lambda_i I+N_i$.
Then $N_i\ne0$ and $N_i^2=0$
by (F1R34.12). Its image is
isotropic, since
$\Omega_i(N_i x,N_i y)=\Omega_i(x,N_i^2y)=0$,
so $\operatorname{rank}N_i\le2$.
The matrix $J_4N_i$ is alternating
by self-adjointness and has the
same rank as $N_i$. The rank of
an alternating matrix is even:
successive nonzero pairings split
off invertible two-by-two alternating
blocks. Thus
\[
                      N_i^2=0,\qquad \operatorname{rank}N_i=2.
 \tag{F1R34.14}
\]
In particular $\operatorname{im}N_i=
\ker N_i$. Choose a basis $e_1,e_2$
of that image and preimages $f_1,f_2$
with $N_i f_j=e_j$. These four
vectors are independent and give
two actual two-dimensional Jordan
summands. This construction includes
the repeated root zero.

\paragraph{3. Every tensor case gives the required nullity bound.}
We prove
\[
       \dim\ker(\mathcal P-uI)
          +\dim\ker(\mathcal P-vI)\le12
 \tag{F1R34.15}
\]
for the actual operators above.
All decompositions below are induced
by (F1R34.13) or the explicit
basis of (F1R34.14).

When both $P_i$ have distinct roots,
the tensor space is the direct sum
of four four-dimensional eigenspaces,
with eigenvalues
\[
       \lambda_1\lambda_2,\quad\lambda_1\mu_2,
                     \mu_1\lambda_2,\quad\mu_1\mu_2.
 \tag{F1R34.16}
\]
At least one of these four values
is outside $\{u,v\}$. Indeed, if
all belonged to that set, all
factors would be nonzero. After
interchanging $u,v$, the upper-left
entry of the two-by-two product
array would be $u$. Since both
row and column factors vary, its
other two adjacent entries would
have to be $v$. The final entry
would then be $v^2/u$. Equality
to $v$ gives $u=v$, and equality
to $u$ gives $u=-v$ or $u=v$.
Each is excluded by (F1R34.8).
Thus at least one entire
four-dimensional summand contributes
to neither kernel, proving
(F1R34.15) in this case.

Next suppose only $P_1$ has a
repeated root, so $P_1=\lambda I+N$
with (F1R34.14), while $P_2$ has
two eigenspaces of dimension two.
On the eight-dimensional summand
corresponding to an eigenvalue
$\eta$ of $P_2$, the tensor operator
is
\[
                  \eta(\lambda I+N)\otimes I_2.
 \tag{F1R34.17}
\]
If $\eta=0$ or $\lambda=0$, its
only eigenvalue is zero and it
contributes to neither nonzero
exceptional kernel. Otherwise its
sole eigenvalue is $\lambda\eta$.
At most one of $u,v$ equals that
value, and its eigenspace is
$\ker N\otimes\mathbf C^2$ of
dimension four. Summing over the
two eight-dimensional summands
gives a total of at most eight.
The case with the repeated-root
factor in the other position is
identical after the actual factor
interchange isomorphism.

Finally suppose both roots are
repeated, so $P_1=\lambda I+N$
and $P_2=\eta I+M$, with both
nilpotents as in (F1R34.14).
If $\lambda\eta=0$, the tensor
operator has sole eigenvalue zero,
as is visible on the explicit
upper triangular Jordan bases.
Neither exceptional kernel occurs.
If $\lambda\eta\ne0$, rescale the
two nonzero nilpotents by their
eigenvalues and choose the
preimages in (F1R34.14) accordingly.
On each of the four actual tensor
products of two Jordan planes,
\[
 \mathcal P=\lambda\eta(I+N_0)\otimes(I+M_0),
 \qquad N_0 f=e,\quad M_0 f'=e',\quad
                N_0e=M_0e'=0.
 \tag{F1R34.18}
\]
After subtracting its sole eigenvalue
and dividing by $\lambda\eta$,
the map on the indicated basis is
\[
 \begin{array}{c|c}
 ee'&0\\
 ef'&ee'\\
 fe'&ee'\\
 ff'&ef'+fe'+ee'.
 \end{array}
 \tag{F1R34.19}
\]
Its two displayed nonzero image
vectors are independent, so it
has rank two and kernel dimension
two. At most one exceptional value
can equal $\lambda\eta$, and its
total eigenspace dimension is eight.
This proves (F1R34.15) in the final
case, including all degeneracies.

Equations (F1R34.7) and
(F1R34.15) imply that the two
mixed-character ranks sum to at
least $32-12=20$. The two even
blocks have rank four for each
endpoint by the exact pairing
calculation (F1R33.9), since
$a_0,a_0+1,c$ are nonzero.
Consequently
\[
          \operatorname{rank}c^2_{z_A}
             +\operatorname{rank}c^2_{z_T}\ge28,
 \qquad
          \operatorname{ext}^2_G(A,A)
             +\operatorname{ext}^2_G(T,T)\ge28.
 \tag{F1R34.20}
\]
The final inequality uses the same
invariant evaluation/full-trace
factorizations as (F1R33.5).

\paragraph{4. The three-map triangle retains every Yoneda kernel.}
Put $W=\operatorname{Hom}_G(A,T)$,
$\dim W=3$, and retain the actual
maps
\[
 \nu:W\to\operatorname{Ext}^2_G(A,A),\quad\nu(f)=\delta f,
 \qquad
 \mu:W\to\operatorname{Ext}^2_G(T,T),\quad\mu(f)=f[2]\delta.
 \tag{F1R34.21}
\]
Neither map is assumed injective.
Set $r_\nu=\operatorname{rank}\nu$,
$r_\mu=\operatorname{rank}\mu$, and
\[
 J_\delta(f)=\delta[2]f[2]\delta,
            \qquad r_J=\operatorname{rank}J_\delta.
 \tag{F1R34.22}
\]
The target endpoint Yoneda membership
supplies, for the two block choices
$k_i\in K_i$ with $D_{k_i}p_i=
\omega_i\ne0$, actual $f_i\in W$
such that
\[
                        \mu(f_i)=\operatorname{ev}_T^2(k_i).
 \tag{F1R34.23}
\]
No uniqueness, shared endpoint
coefficient, or linear choice over
the whole kernel is required here.
The same full Hochschild module and
Serre-trace computation as
(F1R33.16)--(F1R33.19) gives on
these two vectors the pairing
matrix with zero diagonal and
nonzero off-diagonal
$c\omega_1\omega_2$. It follows
that $f_1,f_2$ are independent
and that
\[
                              r_J\ge2.
 \tag{F1R34.24}
\]
This conclusion uses the same
actual coupling $\delta$; it does
not identify all of $W$ with a
two-dimensional characteristic
image.

Use again the vector space
$\mathcal M$ of actual triangle
morphisms (F1R33.23). Its projection
to the endpoint pair is onto the
kernel of
\[
 D(a,t)=a[2]\delta-\delta[2]t,
 \qquad
 D:\operatorname{Ext}^2_G(A,A)\oplus
       \operatorname{Ext}^2_G(T,T)
                       \longrightarrow\operatorname{Ext}^4_G(T,A).
 \tag{F1R34.25}
\]
The proof is the same actual
twice-rotated TR3 diagram, including
both cancelling minus signs; it
does not use injectivity of
$\nu$ or $\mu$.
Invariant CY4 Serre duality gives
$\operatorname{Ext}^4_G(T,A)\simeq W^*$,
so $\operatorname{rank}D\le3$ and
\[
 \dim\mathcal M\ge\dim\ker D
       \ge\operatorname{ext}^2_G(A,A)
            +\operatorname{ext}^2_G(T,T)-3.
 \tag{F1R34.26}
\]

The kernel of the projection
$\mathcal M\to\operatorname{Ext}^2_G(E,E)$
consists exactly of pairs
\[
           (a,t)\in\operatorname{im}\nu\oplus\operatorname{im}\mu,
                                  \qquad D(a,t)=0,
 \tag{F1R34.27}
\]
by the two precise Hom sequences
and preceding groups in
(F1R33.25). For actual representatives
$a=\nu(f)$, $t=\mu(g)$, their
remaining equation is
\[
                            D(\nu(f),\mu(g))=J_\delta(f-g).
 \tag{F1R34.28}
\]
This is independent of the chosen
representatives: both $\ker\nu$
and $\ker\mu$ lie in $\ker J_\delta$
by their displayed compositions.
The restriction of $D$ to
$\operatorname{im}\nu\oplus\operatorname{im}\mu$
has image exactly $\operatorname{im}J_\delta$.
Therefore the kernel in
(F1R34.27) has the exact dimension
\[
 \begin{split}
 r_\nu+r_\mu-r_J
       &=2h-\dim\ker\nu-\dim\ker\mu-r_J\\
       &\le6-2=4.
 \end{split}
 \tag{F1R34.29}
\]
Combining the two actual projections
gives the safe estimate
\[
 \begin{split}
 \operatorname{ext}^2_G(E,E)
    &\ge \dim\mathcal M-(r_\nu+r_\mu-r_J)\\
    &\ge \operatorname{ext}^2_G(A,A)
                  +\operatorname{ext}^2_G(T,T)-7.
 \end{split}
 \tag{F1R34.30}
\]
In particular a smaller loss of
$2h=6$ has not been presumed.
The exact kernels in (F1R34.29)
allow loss seven when both Yoneda
maps are injective and $r_J=2$.

\paragraph{5. The closed rank-$(3,2)$ gate and published comparisons.}
The same actual balanced weakly
semiregular lift has
$\operatorname{ext}^2_G(E,E)=18$
by its identity-sector and
Serre-transpose maps, as detailed
in (F1R33.29). Equations
(F1R34.20) and (F1R34.30) give
at least twenty-one, a contradiction.
This excludes $h=3$. R19 excludes
$h\le1$ at this gate, while
R21--R23 reduce its remaining
factorizations to $h=2$ or $3$.
Together with R33, all actual
rank-$(3,2)$ equal-slope semistable
TF adjacent objects in the stated
gate are therefore excluded.

The proof is about the actual
original coupling and endpoint
Ext spaces. It permits arbitrary
surface and lower-dimensional
coherent defects and every given
finite point lift in this scope.
It constructs no replacement object
and leaves the broader RM seed
problem outside these endpoint
ranks and hypotheses unresolved.

The trace, graded-centre signs,
full-character representation,
Serre duality and triangle-map
framework are the precise published
and local maps audited in R33.
The new four-dimensional operator
identity, the eigenprojectors,
the repeated-root image and
preimage bases, and all tensor
kernel counts are proved explicitly
in (F1R34.9)--(F1R34.19).
Their classical matrix precedent
is acknowledged: V.~Ramakrishna,
Y.~Ansari and F.~M.~Costa,
\emph{The Minimal Polynomial of
Some Matrices Via Quaternions},
Adv.~Appl.~Clifford Algebras 22
(2012), 159--183,
DOI 10.1007/s00006-011-0294-4,
Theorem~3.4 and Proposition~4.1
in the primary author version.
The complete latter proof was
read, as well as the former
statement. That source treats
real four-by-four matrices.
The complex operator identity
and its exact tensor consequences
needed here are the direct
calculations above, so no unstated
extension of that source is a
premise and no novelty is claimed
for the classical quadratic
matrix phenomenon.
